\documentclass[12pt]{article}
\usepackage{ficha}

\begin{document}

\ficha{Lógica e Métodos de Demonstração}[Proposições, Operações Lógicas, Quantificadores e Demonstrações]

\instrucoes{Resolve os exercícios, justificando as respostas. Usa $\wedge$ (e), $\vee$ (ou), $\sim$ (negação), $\Rightarrow$ (implicação) e $\Leftrightarrow$ (equivalência).}

\secao{Proposições}{Indica se cada expressão é uma proposição e, em caso afirmativo, o seu valor lógico.}
\begin{exercicios}
  \item $7$ é um número primo.
  \item $2 + 3 = 6$.
  \item $x + 1 = 4$.
  \item Lisboa é a capital de Portugal.
  \item Que horas são?
  \item $\sqrt{2}$ é um número racional.
  \item Todos os quadrados são retângulos.
  \item $\pi > 3$.
\end{exercicios}

\secao{Operações Lógicas}{Sejam $p$: «$4$ é par», $q$: «$4 > 10$» e $r$: «$9$ é um quadrado perfeito». Indica o valor lógico de cada proposição.}
\begin{exercicios}
  \item $p \wedge q$
  \item $p \vee q$
  \item $\sim q$
  \item $q \Rightarrow r$
  \item $p \Rightarrow q$
  \item $p \Leftrightarrow r$
  \item $(p \wedge r) \Rightarrow q$
  \item $\sim (p \vee q)$
  \item $(\sim p) \vee (q \wedge r)$
  \item Traduz em linguagem corrente: $p \wedge (\sim q)$.
\end{exercicios}

\secao{Tabelas de Verdade e Leis Lógicas}{Constrói as tabelas de verdade e verifica as equivalências.}
\begin{exercicios}
  \item Constrói a tabela de verdade de $p \Rightarrow q$.
  \item Mostra que $p \Rightarrow q$ é equivalente a $(\sim p) \vee q$.
  \item Mostra que $\sim (p \wedge q) \Leftrightarrow (\sim p) \vee (\sim q)$ (primeira lei de De Morgan).
  \item Mostra que $\sim (p \vee q) \Leftrightarrow (\sim p) \wedge (\sim q)$ (segunda lei de De Morgan).
  \item Mostra que $p \Rightarrow q$ é equivalente a $\sim q \Rightarrow \sim p$ (contrarrecíproco).
  \item Mostra que $p \vee (\sim p)$ é uma tautologia.
  \item Escreve a negação de $p \Rightarrow q$ sem usar o símbolo $\Rightarrow$.
  \item Simplifica $\sim (\sim p \wedge q)$.
\end{exercicios}

\secao{Condições e Quantificadores}{Indica o valor lógico de cada proposição, em $\R$, e escreve a sua negação.}
\begin{exercicios}
  \item $\forall x,\ x^2 \ge 0$
  \item $\exists x : x^2 = -1$
  \item $\forall x,\ x + 1 > x$
  \item $\exists x : 2x + 3 = 0$
  \item $\forall x,\ x^2 > x$
  \item $\exists x : x^2 - 5x + 6 = 0$
  \item $\forall x,\ (x > 2 \Rightarrow x^2 > 4)$
  \item $\forall x,\ (x^2 > 4 \Rightarrow x > 2)$
  \item Indica um contraexemplo para a proposição falsa do exercício 8.
  \item Escreve em linguagem simbólica: «Existe um número real cujo quadrado é igual ao seu dobro.»
\end{exercicios}

\secao{Implicação e Equivalência}{Indica se a implicação ou equivalência é verdadeira, em $\R$, justificando.}
\begin{exercicios}
  \item $x = 3 \Rightarrow x^2 = 9$
  \item $x^2 = 9 \Rightarrow x = 3$
  \item $x^2 = 9 \Leftrightarrow (x = 3 \vee x = -3)$
  \item $x > 0 \wedge y > 0 \Rightarrow xy > 0$
  \item $xy > 0 \Rightarrow x > 0 \wedge y > 0$
  \item Indica uma condição necessária mas não suficiente para que um número natural seja múltiplo de 6.
  \item Indica uma condição suficiente mas não necessária para que um número seja par.
\end{exercicios}

\secao{Métodos de Demonstração}{Demonstra cada afirmação, usando o método indicado.}
\begin{exercicios}
  \item (Demonstração direta) A soma de dois números pares é par.
  \item (Demonstração direta) O quadrado de um número ímpar é ímpar.
  \item (Demonstração direta) Para todo $n \in \N$, $n^2 + n$ é par.
  \item (Contrarrecíproco) Se $n^2$ é par, então $n$ é par.
  \item (Redução ao absurdo) $\sqrt{2}$ é irracional.
  \item (Contraexemplo) Mostra que é falso: «Para todo $n \in \N$, $n^2 + n + 41$ é primo.» (Experimenta $n = 40$.)
\end{exercicios}

\secao{Desafios}{Resolve os problemas.}
\begin{exercicios}
  \item Prova que, para todo o número natural $n$, $n^3 - n$ é divisível por 6.
  \item Prova por redução ao absurdo que não existe um maior número natural.
  \item Considera a proposição: «Se um quadrilátero é um quadrado, então tem as diagonais iguais.» Escreve a recíproca, a contrária e a contrarrecíproca. Quais são verdadeiras?
  \item Num reino, os cavaleiros dizem sempre a verdade e os escudeiros mentem sempre. A diz: «B é escudeiro.» B diz: «A e eu somos do mesmo tipo.» O que é cada um?
\end{exercicios}

\end{document}
